% Fonte principal: tcc/markdown/02-revisao-bibliografica.md.
\chapter{REVISÃO BIBLIOGRÁFICA}
\label{chap:revisao-bibliografica}

Este capítulo apresenta os conceitos usados nos capítulos seguintes: o sensoriamento remoto e o comportamento espectral dos materiais, as características do Sentinel-2, a distinção entre litologia e composição química, o aprendizado de máquina supervisionado em dados tabulares, o funcionamento do TabPFN e as formas de avaliar modelos preditivos com dados espaciais.

\section{SENSORIAMENTO REMOTO E COMPORTAMENTO ESPECTRAL}

Meneses, Almeida e Sano \cite{meneses2025} definem o sensoriamento remoto como uma técnica de medição das propriedades espectrais dos objetos da superfície terrestre a partir de imagens obtidas por aviões e satélites. O elo entre o objeto e o sensor é a radiação eletromagnética: parte da energia que incide sobre um material é absorvida, e parte é refletida e registrada pelo sensor. A reflectância é a razão entre a energia que deixa uma área do terreno e a energia que incide sobre ela, medidas no mesmo instante.

Como a absorção depende da composição do material, as faixas do espectro em que a energia é absorvida formam uma assinatura espectral, marcada por feições de absorção \cite{meneses2025}. Segundo Clark \cite{clark1999}, essas feições resultam de processos eletrônicos e vibracionais: os óxidos e hidróxidos de ferro absorvem no visível, enquanto as ligações entre metais e hidroxila, presentes em argilominerais, e os carbonatos produzem absorções no infravermelho de ondas curtas, entre 2,2 e 2,35 $\mu$m. Por isso, Meneses, Almeida e Sano \cite{meneses2025} chamam o infravermelho de ondas curtas de faixa espectral geológica.

A reflectância, porém, descreve apenas a superfície observada. Em solos secos, a informação obtida no infravermelho próximo e de ondas curtas raramente ultrapassa 1 mm de profundidade \cite{norouzi2021}, e coberturas de vegetação a partir de cerca de 10\% podem encobrir o espectro das rochas \cite{grebby2011}. Revestimentos de um material sobre outro também alteram o espectro registrado \cite{clark1999}. Assim, a resposta espectral de um pixel nem sempre corresponde à rocha subjacente.

\section{SENTINEL-2}

A missão Sentinel-2, da Agência Espacial Europeia (ESA), é formada por satélites em órbita heliossíncrona a 786 km de altitude que, operando em conjunto, observam o mesmo ponto do equador a cada cinco dias \cite{drusch2012}. O sensor \textit{Multispectral Instrument} (MSI) registra 13 bandas espectrais, do visível ao infravermelho de ondas curtas, em uma faixa imageada de 290 km: quatro bandas têm resolução espacial de 10 m, seis de 20 m e três de 60 m. As duas bandas do infravermelho de ondas curtas, centradas em 1.610 e 2.190 nm, têm resolução de 20 m \cite{esa2026a}.

As imagens são distribuídas em níveis de processamento. O nível 1C fornece a reflectância no topo da atmosfera, com correções radiométricas e geométricas; o nível 2A aplica a correção atmosférica e a classificação da cena e fornece a reflectância de superfície \cite{esa2026b}. Como a reflectância de superfície reduz o efeito da atmosfera, ela é a medida mais próxima do comportamento espectral dos materiais descrito na seção anterior.

van der Meer, van der Werff e van Ruitenbeek \cite{vandermeer2014} avaliaram o potencial do Sentinel-2 para a geologia com dados simulados a partir de um sensor hiperespectral aerotransportado. As razões entre bandas permitiram mapear óxidos de ferro e minerais com ferro ferroso, com resultados semelhantes aos do \textit{Advanced Spaceborne Thermal Emission and Reflection Radiometer} (ASTER). Por outro lado, as razões usadas com o ASTER para carbonatos, micas e argilominerais não têm equivalente no Sentinel-2, porque três bandas do infravermelho de ondas curtas do ASTER estão contidas em uma única banda do Sentinel-2, que também não possui bandas no infravermelho termal.

\section{LITOLOGIA E COMPOSIÇÃO QUÍMICA}

O termo litologia designa o estudo e a descrição de uma rocha e é frequentemente usado para nomear a própria rocha \cite{winge2018a}. Na classificação litológica, cada amostra recebe uma classe que indica o tipo de rocha, definida por observação de campo, por análise de laboratório ou por um mapa geológico.

A geoquímica estuda a composição química dos constituintes do planeta \cite{winge2018b}. Os dados geoquímicos são organizados em óxidos de elementos maiores, elementos-traço e isótopos, obtidos por técnicas como a fluorescência de raios X e a espectrometria de massa com plasma indutivamente acoplado \cite{rollinson2021}.

A composição química, expressa pelos teores de elementos ou óxidos, difere da composição mineralógica, que indica quais minerais formam a rocha e em que proporção: um mesmo elemento pode estar distribuído em minerais diferentes. Como as feições de absorção são diagnósticas sobretudo da mineralogia \cite{clark1999}, a relação entre o espectro e o teor de um elemento depende dos minerais em que esse elemento ocorre e de sua exposição na superfície.

\section{APRENDIZADO DE MÁQUINA SUPERVISIONADO EM DADOS TABULARES}

No aprendizado supervisionado, o modelo aprende a fazer predições a partir de dados em que os rótulos, isto é, os valores da variável resposta, são observados \cite{izbicki2020}. Quando a resposta é quantitativa, o problema é de regressão; quando é qualitativa, de classificação. Neste trabalho, a classe litológica corresponde a uma tarefa de classificação, e o teor de um elemento ou óxido, a uma tarefa de regressão. Os dados são tabulares: cada linha é uma amostra, e as colunas são atributos, como a reflectância em cada banda.

O desempenho no conjunto usado para ajustar o modelo é uma estimativa otimista do erro em novos dados, por isso parte dos dados é reservada para teste \cite{izbicki2020,james2013}. Quando o modelo apresenta erro pequeno no treinamento e grande no teste, ocorre o \textit{overfitting}: o modelo se ajusta à amostra específica, mas generaliza mal \cite{james2013}.

Entre os modelos usados em dados tabulares, destacam-se os baseados em árvores de decisão. O \textit{Random Forest} combina várias árvores, cada uma dependente de um sorteio aleatório feito de forma independente, e define a classe pelo voto das árvores, podendo também ser usado em regressão \cite{breiman2001}. No \textit{gradient boosting}, as árvores são adicionadas uma a uma, cada nova árvore prevendo os resíduos das anteriores \cite{izbicki2020}; o \textit{Extreme Gradient Boosting} (XGBoost) acrescenta a esse processo uma regularização para reduzir o \textit{overfitting} \cite{chen2016}. Em um conjunto de 45 bases de referência, Grinsztajn, Oyallon e Varoquaux \cite{grinsztajn2022} mostraram que os modelos baseados em árvores continuavam superiores às redes neurais profundas em dados tabulares de tamanho médio, com cerca de 10 mil amostras.

\section{TABPFN}

O TabPFN é baseado no \textit{Transformer}, arquitetura de redes neurais construída sobre mecanismos de atenção, que calculam cada saída como uma soma ponderada de representações da entrada \cite{vaswani2017}. Müller et al. \cite{muller2022} propuseram as \textit{Prior-data Fitted Networks}, redes treinadas com tarefas de aprendizado supervisionado sorteadas de uma distribuição a priori. Em cada passo do treinamento, a rede recebe um conjunto de exemplos rotulados de uma dessas tarefas e aprende a prever o rótulo de outro exemplo. Ao final, ela aproxima a inferência bayesiana para essa distribuição e faz previsões probabilísticas para novos dados em uma única passagem pela rede. Esse procedimento é chamado de \textit{in-context learning}: os exemplos de treinamento são fornecidos na entrada, sem atualização dos parâmetros.

Hollmann et al. \cite{hollmann2023} aplicaram essa ideia aos dados tabulares no TabPFN, com uma distribuição a priori baseada em modelos causais estruturais, que geram conjuntos sintéticos com diferentes relações entre as variáveis. A primeira versão realizava apenas classificação, com até 1.000 amostras de treinamento, 100 atributos numéricos e 10 classes. A versão de Hollmann et al. \cite{hollmann2025} foi pré-treinada com milhões de conjuntos sintéticos e passou a realizar também regressão, prevendo uma distribuição de valores para a variável resposta. Ao receber um novo conjunto, o modelo usa as amostras de treinamento rotuladas e as de teste como entrada e realiza o treinamento e a predição em uma única passagem pela rede, sem ajustar seus parâmetros a esse conjunto.

Nas bases avaliadas por Hollmann et al. \cite{hollmann2025}, com até 10 mil amostras e 500 atributos, o TabPFN superou os métodos comparados, entre eles \textit{Random Forest}, XGBoost, CatBoost e LightGBM. Na classificação, em 2,8 segundos, ele superou um conjunto dos melhores modelos ajustados durante quatro horas. Os autores apontam como limitações a inferência mais lenta que a de modelos otimizados, como o CatBoost, o uso de memória proporcional ao tamanho do conjunto e a avaliação restrita a esse tamanho de dados.

Versões posteriores ampliaram esses limites. O TabPFN-3.5, versão a ser utilizada neste trabalho, tem limites recomendados de 1 milhão de amostras e 6 mil atributos e usa um mesmo modelo para classificação e regressão \cite{priorlabs2026}. Os autores também o avaliaram em conjuntos cujo teste contém grupos ausentes do treinamento, como um novo país ou hospital: nos conjuntos com até 100 mil amostras, seu desempenho igualou o dos melhores métodos comparados. Os pesos do modelo são distribuídos com uma licença que permite o uso acadêmico e em pesquisa, mas não o uso comercial.

\section{AVALIAÇÃO DE MODELOS PREDITIVOS}

Na classificação, a avaliação parte da matriz de confusão, que cruza as classes previstas com as observadas nos dados de referência \cite{foody2002}. Dela se obtém a acurácia global, a proporção de amostras classificadas corretamente. Essa medida depende da abundância das classes: uma classificação pode ter acurácia alta ao acertar as classes dominantes e errar as raras \cite{foody2020}. Por isso, recomenda-se apresentá-la com a matriz de confusão e medidas por classe \cite{olofsson2014}. O coeficiente kappa, embora comum, mede a concordância além do acaso, e não a acurácia, e seu uso rotineiro não é recomendado \cite{foody2020}. O F1 combina, para cada classe, a proporção de acertos entre as amostras atribuídas a ela e entre as que de fato pertencem a ela. Para resumir o resultado de várias classes, a média macro dá o mesmo peso a todas, enquanto a média micro favorece as classes maiores \cite{sokolova2009}. O F1 macro, porém, é calculado de duas formas na literatura: como a média dos F1 de cada classe ou combinando as médias das duas proporções de acertos. Os resultados podem diferir em até 0,5 e mudar a ordem dos modelos comparados; por ser mais robusta quando os erros se concentram em algumas classes, a primeira forma é a recomendada \cite{opitz2021}.

Na regressão, a raiz do erro quadrático médio (RMSE) e o erro absoluto médio (MAE) resumem a diferença entre os valores previstos e os observados. O RMSE dá mais peso aos erros grandes \cite{chai2014}, enquanto o MAE representa diretamente o erro médio \cite{willmott2005}. Chai e Draxler \cite{chai2014} recomendam avaliar os modelos com uma combinação de medidas, e não com uma só. O coeficiente de determinação ($R^2$) indica a proporção da variância da variável resposta explicada pelo modelo e pode ser negativo quando o erro do modelo supera a variância dos dados, isto é, quando prever a média seria melhor \cite{chicco2021}.

Em dados espaciais, a forma de separar treinamento e teste também afeta o resultado. Segundo a primeira lei da geografia de Tobler \cite[p.~236]{tobler1970}, todas as coisas estão relacionadas, mas coisas próximas estão mais relacionadas que coisas distantes. Quando amostras de teste estão próximas das de treinamento, a validação cruzada com partições aleatórias superestima o desempenho em novas áreas; Roberts et al. \cite{roberts2017} recomendam a validação cruzada espacial (\textit{spatial cross-validation}), em que as partições são formadas por blocos espaciais ou regiões distintas. Em um mapeamento de biomassa florestal, Ploton et al. \cite{ploton2020} obtiveram $R^2$ de 0,53 com validação aleatória e de 0,14 com validação cruzada espacial. Além disso, as previsões perdem sentido quando os atributos de uma região diferem muito dos dados de treinamento; Meyer e Pebesma \cite{meyer2021} propõem delimitar a área de aplicabilidade (\textit{area of applicability}), a área em que o modelo aprendeu as relações a partir dos dados de treinamento e em que o desempenho estimado na validação se mantém.
